\section{Transporting marked boundary lifts}\label{sec:transport}

The geometric route needs to retain the way compressed pieces meet one
another.  Knowing the genus and orientability of a piece is insufficient:
different boundary circles may carry different attachments.  The existing
\code{surface_cover.py} computes component families for a restricted
dihedral covering presentation and compares ordered points over a common
basepoint.  We extend that interface to ordered \emph{unparameterized
boundary circles}, optionally together with points.  The extension returns
all compatible covering isomorphisms in a constant number of binary
records.  In the cyclic case it also supplies a canonical key, an exact
count of marking types, and a bound on this count in terms of Euler
characteristic and the number of attachments.

\status{Scope.} The input already includes a checked covering presentation.
All isomorphisms are over the identity of the same base surface, with the
same chosen peripheral paths.  A marked boundary circle is preserved as
a set; its parametrization and a phase on that circle are not prescribed.
An arbitrary attaching map, an embedding into a three-manifold, and a
change of the base presentation carry additional information.  None is
certified by this interface.  In particular, this section does not extract
these covers from a normal surface or implement a hierarchy algorithm.
If an application prescribes orientations of covering components or
directions of interval fibres, those are additional constraints on the
returned maps; the query here asks for all covering isomorphisms.  This
distinction matters for orientable covers of nonorientable bases.

The correspondence between connected covers and transitive monodromy
actions, and the interpretation of covering isomorphisms as equivariant
maps, are classical \cite[Section 1.3]{hatcher-covering}.  Polynomial
compressed normal-surface algorithms are also established
\cite{aht-orbits}; Lackenby's Theorem~14.4 already gives a compressed
cutting procedure with parallelity-bundle information
\cite{lackenby2026}.  Our contribution is a concrete extension of the
repository's restricted covering kernel, its exact arithmetic analysis,
and the consequences for attachment data.  We make no priority claim for
the general covering-space facts used in the proofs.

\subsection{The checked covering model}

Let $F$ be a connected compact surface with nonempty boundary.  Its free
fundamental group has rank $r=1-\chi(F)$.  A canonical surface schema
specifies the free generators and the peripheral words; arbitrary supplied
words are not accepted as evidence for that schema.  The model records
orientability, genus or crosscap number, and the number of boundary
components.  On a fibre $\Z/W\Z$, each generator acts by
\[
 x\longmapsto \epsilon_i x+a_i\pmod W,
 \qquad \epsilon_i\in\{1,-1\}.
\]
The positive integer $W$ and the shifts are binary integers.  Since the
fundamental group is free, every such assignment gives a cover.  An affine
sign is not a surface-orientation character; the existing input validator
keeps these notions separate.

If all signs are positive, put $d=\gcd(W,a_1,\ldots,a_r)$.  Otherwise fix
one negative generator $R(x)=-x+a_0$ and let $d$ be the gcd of $W$, the
shifts of the positive generators, and the differences $a_i-a_0$ for all
negative generators.  Set $m=W/d$.  In both cases the translation subgroup
is generated by $T(x)=x+d$.  Indeed, products of two reflections yield the
listed differences, reflection conjugates a translation to its inverse,
and a Bezout identity gives the converse inclusion.

Without a reflection, each component is one residue class modulo $d$.
With a reflection, components correspond to orbits of
$r_0\mapsto a_0-r_0$ on $\Z/d\Z$.  They therefore consist of either one
fixed residue or a pair of residues.  A single-residue component has $m$
sheets and a paired component has $2m$ sheets.  All labels below retain
the actual residue classes even when two components are isomorphic.

For a peripheral monodromy $p_b$, let $O_b(z)$ denote the orbit of a sheet
$z$.  This orbit identifies a lifted boundary circle over the $b$th base
boundary.  Explicitly,
\begin{align}
 p_b(z)=z+c &\quad\Longrightarrow\quad
 O_b(z)=\{w:w\equiv z\pmod{\gcd(W,c)}\},\label{eq:bt-translation-orbit}\\
 p_b(z)=-z+c &\quad\Longrightarrow\quad
 O_b(z)=\{z,c-z\pmod W\}.\label{eq:bt-reflection-orbit}
\end{align}
The second set has one or two elements.  These are descriptions of
peripheral orbits, not lists of all sheets in a component.

\begin{definition}[Attachment constraints]
Fix source and target components $C,C'$ in one covering presentation.
A point constraint $(x,y)$ requires an equivariant isomorphism $f:C\to C'$
to satisfy $f(x)=y$.  A boundary constraint $(b,x,y)$ requires
$f(O_b(x))=O_b(y)$.  All source sheets lie in $C$, all target sheets lie in
$C'$, and the base boundary index is the same on the two sides.  Each
list of constraints is ordered, and repetitions are allowed.  A sheet
used to select $C$ or $C'$ is not itself an additional point constraint.
\end{definition}

The distinction matters even for an annulus cover.  If its generator acts
as $x\mapsto x+1$ on $W$ sheets, each base boundary has a single lift.
Marking those whole circles leaves all $W$ deck transformations available.
Prescribing the image of one point leaves only one transformation.

\subsection{All maps from their possible root images}

Choose as source root $x_0$ the least residue in the source component.
For a target root image $y$, write $f_y$ for the putative equivariant map
with $f_y(x_0)=y$.

\begin{lemma}[Rooted transport]\label{lem:bt-rooted}
Components of different residue cardinalities admit no covering
isomorphism.  For components of the same kind, the possibilities are as
follows.
\begin{enumerate}
\item In a pure translation cover, every target sheet $y$ is admissible and
$f_y(z)=z-x_0+y\pmod W$.
\item Between paired components, every target sheet is admissible and
\[
 f_y(z)=
 \begin{cases}
 z-x_0+y,&z\equiv x_0\pmod d,\\
 z+x_0-y,&z\equiv a_0-x_0\pmod d,
 \end{cases}\pmod W.
\]
\item Between reflection-fixed components, $y$ is admissible exactly when
$2(y-x_0)=0\pmod W$; the map is then $f_y(z)=z-x_0+y$.
\end{enumerate}
In every admissible case the map is unique.
\end{lemma}
\begin{proof}
Equivariance and transitivity force $f(gx_0)=gy$ for every monodromy
element $g$, so the image of the root determines at most one map.
Translations act regularly on a residue class, giving the first formula.
In a paired component every sheet has a unique expression either
$T^j x_0$ or $T^jRx_0$, with $j\in\Z/m\Z$.  Sending these respectively to
$T^jy$ and $T^jRy$ gives the displayed formulas and a bijection.  They
commute with both $T$ and $R$, hence with the whole monodromy group.

On a fixed component transitivity of $T$ still forces the translation
formula.  Commuting with $R$ requires
\[
 -z+a_0-x_0+y=-z+x_0-y+a_0\pmod W,
\]
which is exactly $2(y-x_0)=0$.  Conversely that condition makes the
translation commute with every generator, so it suffices.  The component
degrees $m$ and $2m$ exclude an isomorphism between the remaining cases.
All arguments include $m=1$ and $W=1,2$, where distinct formal affine
signs can describe the same permutation.
\end{proof}

There are at most two residues $t$ in the target component.  Within one
such residue, write
\begin{equation}\label{eq:bt-root-parameter}
 y=t+du,\qquad u\in\Z/m\Z.
\end{equation}
For every source sheet $z$, the rooted formulas have the uniform form
\begin{equation}\label{eq:bt-uniform-map}
 f_y(z)=\sigma_z(t+du)+k_z\pmod W,
\end{equation}
where $\sigma_z\in\{1,-1\}$ and $k_z$ are determined by $z$ and $x_0$.
The negative sign occurs precisely on the other residue of a paired
source component.  Consequently the unknown in a transport query is one
binary residue $u$, with at most two choices of the outer residue $t$.

\begin{theorem}[Exact boundary transport in two records]\label{thm:bt-two-records}
For $k$ point and boundary constraints in the checked dihedral covering
model, all compatible isomorphisms can be computed without expanding a
sheet, peripheral orbit, or connected component.  The answer consists of
at most two disjoint arithmetic progressions
\[
 \{\alpha+j\delta:0\le j<\nu\}\subseteq\{0,\ldots,W-1\}
\]
of root images.  Each root image determines one isomorphism by
Lemma~\ref{lem:bt-rooted}.  Thus the algorithm returns the exact number of
maps, decides their existence, and supplies one when they exist.

After cover preparation, the bit cost is
$O((k+1)B^3)$ conservatively, where $B$ bounds the bit lengths of the
input integers and $\log(k+2)$.  The returned map family uses $O(B)$ bits;
the implementation retains $O((k+1)B)$ bits of query data.
\end{theorem}
\begin{proof}
We describe the complete calculation for one target residue $t$.
For a reflection-fixed source, admissibility begins with
\begin{equation}\label{eq:bt-fixed-congruence}
 2du\equiv 2(x_0-t)\pmod W.
\end{equation}
There is no initial condition in the other two cases.

A point constraint $(z,z')$ gives
\begin{equation}\label{eq:bt-point-congruence}
 \sigma_zdu\equiv z'-\sigma_zt-k_z\pmod W.
\end{equation}
For a translation boundary with peripheral shift $c$, put
$h=\gcd(W,c)$.  Since $f_y$ is equivariant,
$f_y(O_b(z))=O_b(f_y(z))$.  Equation~\eqref{eq:bt-translation-orbit}
therefore turns $(b,z,z')$ into
\begin{equation}\label{eq:bt-boundary-congruence}
 \sigma_zdu\equiv z'-\sigma_zt-k_z\pmod h.
\end{equation}

For completeness, the linear congruence $Au=v\pmod M$ has no solution
unless $g=\gcd(A,M)$ divides $v$.  If it does, all solutions are
\[
 u\equiv (v/g)(A/g)^{-1}\pmod{M/g}.
\]
When $M/g=1$, this means no restriction.  Here every resulting modulus
divides $m$: a point equation has modulus $m$, a translation boundary has
modulus $h/d$, and \eqref{eq:bt-fixed-congruence} has modulus
$m/\gcd(2,m)$.  The statement about $h/d$ uses $d\mid c$, since the
peripheral translation belongs to the monodromy group.

Intersect these congruences using the generalized Chinese remainder
theorem.  Specifically, $u=a\pmod\ell$ and $u=b\pmod q$ are compatible
exactly when $\gcd(\ell,q)\mid b-a$.  In that event they give one residue
class modulo $\lcm(\ell,q)$.  Thus all point, admissibility, and
translation-boundary conditions either fail or leave
\begin{equation}\label{eq:bt-congruence-final}
 u=a\pmod\ell,\qquad \ell\mid m,
 \qquad 0\le a<\ell.
\end{equation}
If there are no reflection-boundary constraints, the possible root images
are exactly
\[
 t+da+jd\ell,\qquad 0\le j<m/\ell,
\]
one progression for this target residue.

Suppose instead that a reflection boundary has shift $c$.  Its constraint
requires
\begin{equation}\label{eq:bt-two-targets}
 \sigma_z(t+du)+k_z\in\{z',c-z'\pmod W\}.
\end{equation}
Solving the point congruence for each of these two targets gives at most
two individual values of $u\pmod m$.  Keep only values satisfying
\eqref{eq:bt-congruence-final}.  Check every remaining reflection-boundary
constraint on this constant-size set.  There is no branching into
$2^k$ combinations: after the first reflection boundary, only two root
values remain to test.

It remains to prove the sharper bound of two output progressions, rather
than four.  In a paired component every reflection boundary exchanges the
two residues.  Its orbit has one sheet in each residue.  For a fixed target
residue $t$, the residue of \eqref{eq:bt-uniform-map} is fixed, so at most
one of the two targets in \eqref{eq:bt-two-targets} can occur.  Each target
residue consequently yields at most one singleton.  In a fixed component,
the initial admissibility condition allows at most two root images in
total.  If both survive, they differ by $W/2$, and form one arithmetic
progression; any subset is empty or a singleton.  A pure translation
presentation has no reflection peripheral monodromy.  These cases exhaust
the model.  Distinct target residues give disjoint root sets.

Every constraint was replaced by an equivalent congruence or a direct
peripheral-orbit membership test.  Lemma~\ref{lem:bt-rooted} provides an
isomorphism for every surviving root, so the procedure is both sound and
complete.  The map count is the sum of the progression lengths.

There are $O(k+1)$ gcd, inverse, and congruence operations on $O(B)$-bit
integers.  Accumulated moduli always divide $m$, so CRT never creates a
modulus with a growing product of unrelated sizes.  Classical Euclidean
and integer-arithmetic bounds give $O(B^3)$ bit work per operation, a
conservative bound sufficient here.  The number of surviving roots can
be of order $W$ but its binary encoding has $O(B)$ bits.  The constant
number of arithmetic-progression records has the same size bound.
\end{proof}

\begin{example}[A constraint on circles is weaker than one on their labels]
In a translation cover with $W=360$ and generator shifts $1,12,18$, choose
the base surface to be a sphere with four boundary circles.  Two prescribed
boundary-lift correspondences give
$u=5\pmod{12}$ and $u=11\pmod{18}$.  They are compatible and leave
\[
 u=29\pmod{36},
\]
so there are ten maps.  Replacing the second target residue by $10$ gives
incompatible congruences modulo their common divisor $6$.  The calculation
does not traverse the $360$-sheet fibre and does not incorrectly require
the two chosen representatives themselves to match.
\end{example}

\subsection{A general bound on the remaining ambiguity}

The small family representation is specific to the dihedral model.  The
following structural restriction holds for arbitrary finite connected
surface covers.

\begin{proposition}[Boundary ambiguity divides peripheral degree]\label{prop:bt-ambiguity}
Let two connected finite covers of the same bordered surface have ordered
specified boundary lifts of degrees $e_1,\ldots,e_k$, with $k\ge1$.
If any covering isomorphism respecting those correspondences exists,
the number of such isomorphisms divides $\gcd(e_1,\ldots,e_k)$.  Their
ambiguity is a cyclic group.  If the image of any one fibre point is also
prescribed, there is at most one such isomorphism.
\end{proposition}
\begin{proof}
Choose one compatible isomorphism $f_0$.  All compatible maps are uniquely
$f_0\circ h$, where $h$ is a deck transformation of the source cover that
preserves each marked boundary lift.  Denote this finite subgroup by $H$.
On a chosen boundary lift of degree $e_i$, its peripheral monodromy is an
$e_i$-cycle.  Every element of $H$ restricts to a permutation commuting
with this cycle.  The centralizer of a full cycle on its own support is
the cyclic rotation group: the image of one point determines the
permutation, and the other images follow around the cycle.

Restriction embeds $H$ in this cyclic group.  Indeed, a deck
transformation in the kernel fixes a point of the fibre.  Equivariance
and connectedness then force it to fix every fibre point, and hence to
be the identity covering map.  Thus $H$ is cyclic and $|H|$ divides each
$e_i$, proving the divisibility assertion.  With a point constraint,
every $h\in H$ fixes that point, so $H$ is trivial.
\end{proof}

For a reflection boundary, $e_i\le2$, explaining the two-value reduction
in the preceding algorithm from a covering-space viewpoint.  The
proposition is also a useful consistency check on implementations: a
nonempty returned family must have a size dividing every marked lift
degree.  It does not imply a bound on the number of different markings.

\subsection{Cyclic canonical keys and an exact state count}

In the pure translation model one can avoid repeated pairwise comparison
by constructing a canonical signature.  Normalize a component with residue
$r_0$ by writing its sheets as $r_0+dj$, $j\in\Z/m\Z$.  A marked point
has a coordinate $v_i\in\Z/m\Z$.  A marked boundary whose peripheral
translation is $c_i$ has a coordinate
\[
 v_i\in\Z/q_i\Z,\qquad q_i=\gcd(W,c_i)/d.
\]
For a point slot put $q_i=m$.  The number $q_i$ counts the boundary lifts
in that component; it is \emph{not} the degree of each such lift, which
is $m/q_i$.  Every $q_i$ divides $m$.

Fix the ordered list of point slots and base boundary labels.  Isomorphisms
act on their coordinates by a common phase:
\begin{equation}\label{eq:bt-diagonal-action}
 (v_1,\ldots,v_k)\longmapsto
 (v_1+u\bmod q_1,\ldots,v_k+u\bmod q_k),\qquad u\in\Z/m\Z.
\end{equation}
This model includes comparison between distinct residue components:
normalization removes their absolute residues, while retaining the same
base presentation.

\begin{theorem}[Canonical cyclic markings and exact orbit count]\label{thm:bt-canonical}
Let $q_1,\ldots,q_k\mid m$, and put
$L=\lcm(q_1,\ldots,q_k)$, with $L=1$ for $k=0$.
There is an exact canonical signature for
\eqref{eq:bt-diagonal-action}, obtained using $O(k+1)$ gcd and CRT
operations on $O(B)$-bit integers.  It supplies an explicit phase taking
the input to its canonical representative.  Every marking has exactly
$m/L$ automorphisms.  The number of distinct marking types is exactly
\begin{equation}\label{eq:bt-exact-type-count}
 \frac{\prod_{i=1}^{k}q_i}{\lcm(q_1,\ldots,q_k)}.
\end{equation}
The empty product is one.  The key must include the fixed base
presentation and the ordered mark schema.
\end{theorem}
\begin{proof}
We compute the lexicographically least tuple in each phase orbit.  Initially
all phases are available: write $u=A\pmod\ell$ with $A=0$, $\ell=1$.
Assume that this congruence describes exactly those phases attaining the
least possible values in the already processed coordinates.  For the next
coordinate of period $q_i$, put $g=\gcd(\ell,q_i)$.  Its attainable values
are the residue class $v_i+A\pmod g$ inside $\{0,\ldots,q_i-1\}$, because
the subgroup generated by $\ell$ modulo $q_i$ consists of the multiples
of $g$.  Its smallest value is
\begin{equation}\label{eq:bt-greedy-coordinate}
 c_i=(v_i+A)\bmod g,\qquad 0\le c_i<g.
\end{equation}
Restrict the phase by the additional congruence
\[
 u=c_i-v_i\pmod{q_i}.
\]
It is compatible with $u=A\pmod\ell$ by construction.  Generalized CRT
replaces the pair by one congruence modulo $\lcm(\ell,q_i)$.

Induction proves that the resulting tuple $(c_1,\ldots,c_k)$ is exactly
the lexicographic minimum, and that the final residue $A\pmod L$
describes all phases attaining it.  Therefore two inputs have equal
canonical tuples exactly when they belong to the same phase orbit.
The phase $A$ explicitly transports the original root to $dA$ in the
residue-zero component.  This also proves that the signature compares
components, not merely tuples of integers with a fortuitous equality.

A phase stabilizes any tuple exactly when it is divisible by every
$q_i$, equivalently when it is zero modulo $L$.  The stabilizer therefore
has $m/L$ elements, and each orbit has size $L$.  There are
$\prod_iq_i$ tuples, giving \eqref{eq:bt-exact-type-count}.  No enumeration
of tuples or phase values is used to calculate this expression.

All phase moduli divide $m$, and every processed coordinate has $O(B)$
bits.  The congruence calculation has the claimed arithmetic size.
The full mark schema and normalized base presentation are necessary:
permuting boundary labels or changing the base cover is outside the
equivalence relation proved here.
\end{proof}

The canonical key and witness phase cost $O((k+1)B^3)$ bit work
conservatively and $O(kB)$ output bits, in addition to the shared model
header.  The implementation also returns the exact integer in
\eqref{eq:bt-exact-type-count}.  That integer can itself require
$O(kB)$ bits.  Accumulating its numerator by ordinary multiplication
adds a conservative $O(k^2B^2)$ bit cost, so the combined bound is
$O((k+1)B^3+k^2B^2)$.  This separates signature construction from writing
a possibly much larger count.

\begin{remark}[Direct indexing of the quotient]\label{rem:bt-indexing}
Put $L_0=1$, $L_i=\lcm(q_1,\ldots,q_i)$ and
$g_i=\gcd(L_{i-1},q_i)$.  The canonical tuples are exactly
\[
 0\le c_i<g_i\qquad(1\le i\le k).
\]
Necessity follows from \eqref{eq:bt-greedy-coordinate}.  Conversely,
running the normalization algorithm on such a tuple leaves its phase
equal to zero at every step, so the tuple is already canonical.  Thus
these are independent mixed-radix digits, and
$\prod_i g_i=\prod_iq_i/L_k$.  They give direct ranking, unranking, or
enumeration of all marking types without inspecting any noncanonical
marking.  A representative in the residue-zero component is obtained
by choosing sheet $dc_i$ for its corresponding point or boundary slot.
The implementation returns the canonical tuple and exact count; this
mixed-radix interpretation requires no additional geometric oracle.
\end{remark}

\subsection{A local quasi-polynomial bound from Euler complexity}

Boundary circles behave differently from arbitrarily placed points.
The distinction gives a useful conditional state bound directly relevant
to hierarchy interfaces.

\begin{corollary}[Few boundary attachments on a bounded surface]\label{cor:bt-genus-bound}
Let $C\to F$ be a connected component of a checked pure translation cover,
and prescribe $k\ge1$ ordered slots for unparameterized boundary circles,
with their base boundary labels fixed.  The number of isomorphism types
of those markings is at most
\begin{equation}\label{eq:bt-euler-bound}
 \max\{2,-\chi(C)\}^{\,k-1}.
\end{equation}
Consequently, for any family with $-\chi(C)\le n^a$ and
$k\le b\log(n+2)$ for fixed constants $a,b$, these normalized local
markings have at most $2^{O((\log(n+2))^2)}$ types.  Their equality keys
are computable with the bit bounds in Theorem~\ref{thm:bt-canonical}.
\end{corollary}
\begin{proof}
Suppose first that $\chi(F)<0$.  The component degree is $m$, and
$\chi(C)=m\chi(F)$, so $m\le-\chi(C)$.  Every mark period satisfies
$q_i\le m$.  Taking an index attaining $\max_iq_i$ gives
\[
 \frac{\prod_iq_i}{\lcm_iq_i}
 \le \frac{\prod_iq_i}{\max_iq_i}
 \le m^{k-1}.
\]
It remains to handle compact connected bordered surfaces with
nonnegative Euler characteristic.  They are the disk, annulus, and
M\"obius band.  A connected disk cover has degree one.  In an annulus
cover, the generator acts transitively on the component and each base
boundary has one lift, so every $q_i=1$.  In a cyclic cover of a
M\"obius band, the boundary monodromy is the square of the generator.
It has $\gcd(2,m)\le2$ orbits in the component, so all mark periods are
at most two; the preceding product-over-maximum argument gives
$2^{k-1}$.  This proves \eqref{eq:bt-euler-bound} in all cases.
Taking logarithms proves the quasi-polynomial bound.
\end{proof}

This corollary concerns only markings of a fixed checked component type
and fixed base boundary labels.  A global state bound would additionally
control the number of underlying presentations, pieces, and labels; it
would also have to prove that these bounds survive corrections and
restarts.  Parameterized attaching maps, twists, and point phases are
not covered by \eqref{eq:bt-euler-bound}.  For example, two ordered fibre
points on a connected $W$-fold annulus cover already have $W$ types,
although the covering surface remains an annulus.  The point periods are
$q_1=q_2=W$, and \eqref{eq:bt-exact-type-count} gives $W$ exactly.

\begin{proposition}[Circle markings can still have exponential diversity]\label{prop:bt-entropy}
For binary size $O(\beta)$ there are checked cyclic covering presentations
with one connected component and $2^{\beta-1}$ inequivalent ordered
two-boundary-circle markings.
\end{proposition}
\begin{proof}
Set $W=2^\beta$, $h=W/2$, with $\beta\ge1$.  Take a sphere with four
boundary components and assign the three canonical free generators
the translations $1,h,h$.  The fourth peripheral monodromy is the
inverse of their product.  Translation by one makes the cover connected.
The boundaries corresponding to the second and third generators each
have $h$ lifted circles, labelled by residues modulo $h$.  A deck
transformation adds the same residue to both labels.  The pairs
$(0,j)$ for $j\in\Z/h\Z$ are therefore inequivalent.  Equivalently,
Theorem~\ref{thm:bt-canonical} gives $h^2/h=h$ types.
The base has fixed size and the three shifts require $O(\beta)$ bits.
\end{proof}

The last example has $\chi(C)=-2W$ and genus $W/2$; its covering surface
has $W+2$ boundary circles.  Thus it does not obstruct the bounded-Euler
hypotheses of Corollary~\ref{cor:bt-genus-bound}.  It shows why a
polynomial-time equality test alone does not prove a small state count.
In this example every marking has only two automorphisms, while the
number of inequivalent markings is exponential.  Map ambiguity, state
diversity, and the bit size of each state are separate quantities.

\subsection{Implementation contract and use in a hierarchy}

The new \code{BoundaryTransportIndex} inherits validation and prepared
monodromy data from \code{CoverIndex}; it does not accept a serialized
classification as a proof of the input. Its \code{transports} method
accepts source and target component identifiers, point pairs, and
boundary triples. It returns a canonical source root,
at most two root progressions, their total cardinality, and one target
anchor or an empty answer.  An inherited \code{transport_sheet} call
evaluates the corresponding map at any requested sheet.  Queries outside
the selected components, malformed indices, and invalid integer fields
raise errors.  An empty family means incompatible constraints within
the checked cover model; it is never interpreted as a knot verdict.

The \code{cyclic_marked_signature} method is deliberately restricted to
presentations with only positive affine signs.  It normalizes point
marks and boundary marks in their fixed ordered schema, returns a
hashable complete key and a transport to that key, and reports both the
number of compatible maps and the number of possible marking types.
It rejects a presentation with a negative generator rather than silently
treating a dihedral cover as cyclic.  The general dihedral comparison
remains available through \code{transports}.

All integer inputs use the repository's exact integer or signed
hexadecimal transport.  The standalone JSON command has cooperative
cancellation, with checks during validation and each constraint step.
A cancellation propagates as an exception; it does not become an
incompatibility declaration.  Individual big-integer operations and
bulk JSON serialization are not hard-interruptible.  The experimental
checks reported separately compare the arithmetic answer against
literal permutation graphs, including nonorientable bases, fixed and
paired components, and redundant affine signs on very small fibres.

There are two immediate ways to use this interface once a suitable
geometric extraction theorem is supplied.  First, a piece comparison can
retain the whole family of permitted boundary transports instead of
choosing or enumerating individual phases.  New attachment evidence
intersects that family by further congruences.  Second, in the cyclic
subcase a cache can use the canonical marked key, so isomorphic markings
share a state.  The resulting local quotient has the exact size in
\eqref{eq:bt-exact-type-count}; under the stated Euler and attachment
bounds it has the quasi-polynomial size in
Corollary~\ref{cor:bt-genus-bound}.  Establishing that actual normal-surface
pieces admit these presentations, transporting their full gluing data,
and maintaining the bounds throughout hierarchy repair remain separate
geometric research tasks.
